\documentclass[../../../main.tex]{subfiles}

\begin{document}

\section{Parameterizations of Curves}

As you could have noticed from the previous section, there are
many different ways to represent a curve. In this section,
we will discuss another method that uses the length of arc of
a curve. But before doing that, let's get started with definitions.

\begin{definition}{}{}
A \textbf{parametrization} of a curve is a way of representing the
curve in a continuous vector-valued function. A parametrization is
\textbf{smooth} if its derivative is nonzero and continuous at the
domain of the curve.
\end{definition}

Notice that a smooth parametrization of a curve is not unique
because we can always scale the parameter with accordingly adjusted
domain or construct with different orientations. Continuing from
the definition, we can establish an important theorem.

\begin{theorem}{}{}
Given a curve $C$, let $\mathbf{r_1}$ be a smooth parametrization of
$C$ with the domain $U$. If $s$ is a differentiable real-valued
function with domain $V$ and range $U$ such that $s'$ is continuous
and nonzero, then a smooth parametrization of $C$ can also be
represented as $\mathbf{r_2} = \mathbf{r_1} \circ s$.
\end{theorem}
\begin{proof}
Without loss of generality, assume $\mathbf{r_1}$ is a
parameterization in 3-space defined as $\mathbf{r_1}(t) =
\langle x(t), y(t), z(t) \rangle$. Because it is self-evident that
$\mathbf{r_2}$ is a parametrization of $C$ by definition, it suffices
to show that $\mathbf{r_2}$ is smooth. Consider the following
equations.
\begin{align*}
    \mathbf{r_2}'(k)
    &= \frac{d}{dk} \langle x(s(k)), y(s(k)), z(s(k)) \rangle \\
    &= s'(k) \langle x'(s(k)), y'(s(k)), z'(s(k)) \rangle \\
    &= s'(k) \mathbf{r_1}'(s(k))
\end{align*}
Notice that because $s'$ and $\mathbf{r_1}'$ are both nonzero and
continuous function by construction, $\mathbf{r_2} = \mathbf{r_1}
\circ s$ is a smooth parametrization of $C$.
\end{proof}

We actually have a special name for $s$. Such functions are known
as a \textbf{smooth change in parameter}. With this concept in mind,
we can find a special way of parametrizing a curve.

Before continuing with the new method, we need to understand how
arc lengths are defined in a curve.

\begin{important}
For a curve defined as $\mathbf{r}(t) = \langle x(t), y(t), z(t)
\rangle$, the \textbf{arc length} $s$ of the curve in the interval
$a \leq t \leq b$ is defined as the following.
\[
s
= \int_a^b \sqrt{x'(t)^2 + y'(t)^2 + z'(t)^2} \,dt
= \int_a^b \| \mathbf{r}'(t) \| \,dt
\]
\end{important}

A way of understanding this equation is cutting the curve in the
given interval into arbitrarily small pieces so that each piece
becomes a linear segment. Then, it becomes obvious that we can
simply use Pythagorean theorem to derive the formula just as we do
for 2-space. Also, using the definition of a norm of a vector,
we can rewrite the integral.

Now with this idea in mind, we can construct arc length
parametrization of a curve. The basic idea of this method is to
define each point on the curve with the distance from a reference
point $t_0 = f(0)$. Consider the following equation.
\[
s = \int_{t_0}^{f(s)} \| \mathbf{r_1}'(t) \| \,dt
\]
From here it is evident that we can define a new $t$-value with some
distance $s$ since there are only two possible directions depending
on the position of $f(s)$. Another fun fact that we can notice by
letting $s = f^{-1}(u)$ for some independent values of $u$ is that
\begin{align*}
    f^{-1}(u)
    &= \int_{t_0}^{f(f^{-1} (u))} \| \mathbf{r_1}'(t) \| \,dt
    = \int_{t_0}^u \| \mathbf{r_1}'(t) \| \,dt \\
    f^{-1}(t)
    &= \int_{t_0}^t \| \mathbf{r_1}'(u) \| \,du
\end{align*}
holds and the following equation is also true by FTC.
\[
\left( f^{-1} \right)'(t) = \| \mathbf{r_1}'(t) \|
\]
This shows that because $f$ is differentiable and $f'$ is continuous
and nonzero from the equation $f'(s) = \frac{1}{{f^{-1}}'(f(s))}$,
$f$ is a smooth change of parameter.

\begin{important}
In other words, the arc length function of a curve with
parametrization $\mathbf{r}$ can be represented as the following
with the reference point $t_0$.
\[
f(t) = \int_{t_0}^t \| \mathbf{r}'(u) \| \,du
\]
\end{important}

This is how we represent a curve with arc length! Before we conclude
this section, here is an important theorem to consider.

\begin{theorem}{}{}
Given an arc length parametrization $\mathbf{r}(s)$ of a curve,
$\| \mathbf{r}'(s) \| = 1$ always holds.
\end{theorem}
\begin{proof}
Note that by definition,
\[
s = \int_0^s \| \mathbf{r}'(u) \| \,du
\]
and we obtain $1 = \| \mathbf{r}'(s) \|$ by differentiating with
respect to $s$.
\end{proof}

This is it for this section and this note! In our next note, we will
discuss about curvature.

\end{document}


